\chapter*{How to Read This Companion}
\addcontentsline{toc}{chapter}{How to Read This Companion}

This volume is the research companion to the short book of
\emph{Adversarial Cooperation}. Both editions follow the same twenty-two
chapters and five appendices. The short book carries their central thoughts;
this companion holds the detailed models, arguments, algorithms, citations,
and implementation references where those have been developed.

Every existing chapter idea remains represented here. Some expansions are
substantial; others are still notes or open questions. Inclusion does not mean
that a chapter is finished, that a claimed property has been established, or
that a contribution is novel.

The \href{Adversarial-Cooperation-Short.pdf}{short book} can be read on its own
or beside this volume. Keep the PDFs together for links between them; matching
chapter numbers and appendix letters also work without electronic links.

The manuscript uses the following status language.

\begin{description}
  \item[Mature foundation.] A literature-backed chapter with explicit limits.
  \item[Educational construction.] A precise teaching artifact whose security
  boundary is narrower than a production protocol.
  \item[Author's intuition.] A philosophical or ethical idea preserved in the
  author's voice; it is not a theorem.
  \item[Design principle.] A proposed way to think about a system or
  institution; evidence may still be absent.
  \item[Open problem or research fragment.] A question, partial formulation,
  or unfinished construction retained for future work.
  \item[Placeholder.] A named subject for which no substantive manuscript has
  yet been written.
\end{description}

The technical spine currently consists of Trust Establishment, Two Oracles
Play Rock--Paper--Scissors, Tic-Tac-Toe Without Revealing the Strategy, Poker
Without Revealing the Cards, and the Hash Functions appendix. Tic-Tac-Toe is a
reusable example for separating a hidden object, a public predicate, temporal
binding, a future private proof, and protocol lifecycle. It is not the subject
of the whole book, and no conclusion transfers automatically to another
chapter.

Three chapters named in the provisional architecture do not yet have source
files: \emph{What Is Adversarial Cooperation?}, \emph{Adversaries, Enemies,
and Partially Aligned Interests}, and \emph{Threat Models and Protocol
Anatomy}. They remain planned work. This draft does not fabricate text to fill
those gaps.

The ethical purpose of the project is peace. Peace is not itself a
cryptographic property. Each technical chapter must therefore distinguish
what its mechanism checks from what remains a human, political, or moral
question.
